\documentclass[a4paper,11pt]{article}
\usepackage[T2A]{fontenc}
\usepackage[utf8]{inputenc}
\usepackage[russian]{babel}
\usepackage{amsmath,amssymb,mathtools}
\renewcommand{\familydefault}{\sfdefault}
\usepackage{icomma}
\usepackage[a4paper,top=22mm,bottom=22mm,left=21mm,right=21mm]{geometry}
\usepackage{microtype}
\usepackage{tikz}
\usetikzlibrary{arrows.meta,positioning,shapes.geometric,calc}
\usepackage{tabularx,booktabs,array}
\usepackage{enumitem}
\usepackage{fancyhdr}
\usepackage{parskip}
\usepackage{xcolor}
\usepackage{hyperref}
\definecolor{accent}{HTML}{008080}
\definecolor{darkaccent}{HTML}{004D4D}
\definecolor{textgray}{HTML}{333333}
\definecolor{muted}{HTML}{647079}
\hypersetup{hidelinks,pdftitle={Чек-лист — Анастасия Павлова — 23.11.26},pdfauthor={Лёвин Артём Александрович}}
\setlength{\parindent}{0pt}
\setlength{\parskip}{5pt}
\setlist[itemize]{leftmargin=16pt,itemsep=3pt,topsep=2pt}
\setlist[enumerate]{leftmargin=19pt,itemsep=5pt,topsep=2pt}
\renewcommand{\arraystretch}{1.22}
\pagestyle{fancy}
\fancyhf{}
\fancyhead[L]{\footnotesize\color{muted} Математика \quad / \quad 23 ноября 2026 г.}
\fancyhead[R]{\footnotesize\color{muted} Анастасия Павлова}
\fancyfoot[C]{\footnotesize\color{muted}\thepage}
\renewcommand{\headrulewidth}{0.25pt}
\setlength{\headheight}{14pt}
\newcommand{\secname}[1]{\section*{\color{darkaccent}#1}\vspace{-4pt}\color{black}}
\newcommand{\tip}[1]{\textcolor{darkaccent}{\textbf{Обратите внимание.}} #1}
\newcommand{\ex}[1]{\textcolor{darkaccent}{\textbf{Пример.}} #1}
\newcommand{\one}{\mathbb{1}}
\begin{document}
\thispagestyle{empty}
\begin{center}
\vspace*{23mm}
{\large\color{darkaccent} МАТЕМАТИКА}\par
\vspace{19mm}
{\LARGE Чек-лист}\par
\vspace{15mm}
{\Large\color{accent} Степени, линейная функция,\\[4pt] векторы и текстовые задачи}\par
\vspace{15mm}
{\normalsize Подготовка к ЕГЭ по профильной математике}\par
\vspace{37mm}
{\large Анастасия Павлова}\par
\vspace{14mm}
{\normalsize 23 ноября 2026 г.}\par
\vfill
{\small Лёвин Артём Александрович}\par
{\small эксклюзивно для Анастасии Павловой}\par
\vspace{15mm}
\begin{tikzpicture}[x=1cm,y=1cm,draw=accent,thick]
\draw (0,0) .. controls (1.4,0.65) and (2.4,-0.5) .. (4.0,0) .. controls (5.3,0.65) and (6.7,-0.55) .. (8.0,0);
\end{tikzpicture}
\vspace{6mm}
\end{center}
\clearpage
\secname{Как работать с чек-листом}
Материал повторяет пять приёмов занятия. После каждого пункта попробуйте воспроизвести правило и один пример самостоятельно. В конце предложены \textbf{10 заданий по пути А} и отдельная таблица ответов.
\begin{enumerate}[label=\textcolor{darkaccent}{\arabic*.}]
\item Приводите степени к общему основанию, контролируйте скобки и знаки.
\item Для прямой $y=kx+b$ используйте координаты точек как два уравнения.
\item Для векторов сначала находите координаты, затем длины и скалярное произведение.
\item В задачах о работе связывайте производительность, время и объём работы.
\item В смесях отдельно учитывайте массу всей смеси и массу растворённого вещества.
\end{enumerate}
\secname{1. Степени, корни и скобки}
Для допустимых значений выражений выполняются свойства:
\[
\begin{aligned}
 a^m\cdot a^n&=a^{m+n},&\frac{a^m}{a^n}&=a^{m-n}\quad(a\ne0),\\
 (a^m)^n&=a^{mn},&(ab)^n&=a^nb^n,\\
 a^{-n}&=\frac1{a^n}\quad(a\ne0),&\sqrt{a}&=a^{1/2}\quad(a\ge0).
\end{aligned}
\]
\tip{При умножении степеней с одинаковым основанием показатели \textbf{складываются}. При возведении степени в степень показатели \textbf{перемножаются}.}

\ex{Решим уравнение, в котором участвуют корень и степень степени:}
\[
9^{x-1}\sqrt{27}=3^{11/2}.
\]
Представим $9=3^2$ и $\sqrt{27}=(3^3)^{1/2}=3^{3/2}$. Тогда
\[
3^{2(x-1)}\cdot 3^{3/2}=3^{11/2}
\quad\Longleftrightarrow\quad 2x-2+\tfrac32=\tfrac{11}{2}
\quad\Longleftrightarrow\quad x=3.
\]

\textbf{Отдельный навык: квадрат двучлена.} Используйте формулу
\[
(a-b)^2=a^2-2ab+b^2,
\qquad (2x-3)^2=4x^2-12x+9.
\]
Число $3$ во второй степени даёт $9$; средний член $-12x$ пропускать нельзя.

\secname{2. Линейная функция по двум точкам}
Если график $f(x)=kx+b$ проходит через точку $(x_0,y_0)$, то $kx_0+b=y_0$. Две точки с разными абсциссами дают систему двух уравнений на $k$ и $b$.

\ex{Пусть $A(-2;7)$ и $B(2;-1)$ лежат на графике функции.} Тогда
\[
\begin{cases}-2k+b=7,\\2k+b=-1.\end{cases}
\]
Вычтем первое уравнение из второго: $4k=-8$, значит $k=-2$. Из $-2k+b=7$ получаем $b=3$ и
\[
\boxed{f(x)=-2x+3},\qquad f(-3)=-2(-3)+3=9.
\]
\tip{Сложение уравнений удобно, если противоположные коэффициенты сокращаются. Вычитание удобно, когда коэффициенты при одной переменной одинаковы. При вычитании заключайте всё вычитаемое уравнение в скобки.}

Чтобы найти \textbf{точку пересечения} $y=f(x)$ с прямой $y=g(x)$, решите $f(x)=g(x)$ и подставьте полученную абсциссу в любую формулу:
\[
-2x+3=x+1\Longrightarrow x=\tfrac23,\qquad y=\tfrac53.
\]
\textbf{Точка пересечения:} $\left(\frac23;\frac53\right)$.
\secname{3. Векторы: координаты, длина, угол}
Пусть $A(x_A;y_A)$ и $B(x_B;y_B)$. Координаты вектора определяются вычитанием \textbf{«конец минус начало»}:
\[
\overrightarrow{AB}=(x_B-x_A;\ y_B-y_A),\qquad
|\overrightarrow{AB}|=\sqrt{(x_B-x_A)^2+(y_B-y_A)^2}.
\]
Для ненулевых векторов $\vec u=(u_x;u_y)$ и $\vec v=(v_x;v_y)$:
\[
\vec u\cdot\vec v=u_xv_x+u_yv_y,
\qquad \cos\varphi=\frac{u_xv_x+u_yv_y}{\sqrt{u_x^2+u_y^2}\sqrt{v_x^2+v_y^2}}.
\]
\tip{Перемножайте \textbf{соответствующие} координаты: $x$ первого с $x$ второго, $y$ первого с $y$ второго. Угол между векторами определяется только для ненулевых векторов.}

\ex{Даны $A(-2;1)$, $B(4;9)$, $C(6;-5)$. Найдём $|AB|$ и $\angle BAC$:}
\[
\overrightarrow{AB}=(6;8),\quad \overrightarrow{AC}=(8;-6),\quad
|AB|=\sqrt{6^2+8^2}=10,\quad |AC|=10.
\]
\[
\cos\angle BAC=\frac{6\cdot8+8\cdot(-6)}{10\cdot10}=0.
\]
Значит, $\boxed{\angle BAC=90^\circ}$. Если скалярное произведение ненулевых векторов равно нулю, эти векторы перпендикулярны. Если $\cos\varphi=-1$, векторы противоположно направлены.

\secname{4. Совместная работа}
Для выполненной работы используем соотношение
\[
\text{объём работы}=\text{производительность}\times\text{время}.
\]
Если весь объём принять за $1$, мастер, работающий самостоятельно $t$ часов, имеет производительность $1/t$ работы в час.

\ex{Первый мастер выполняет заказ на 6 часов быстрее второго. Вместе они завершают его за 4 часа.} Обозначим время первого мастера через $x>0$. Тогда:
\[
\begin{array}{lccc}
\toprule
 &\text{Время, ч}&\text{Производительность}&\text{Объём}\\
\midrule
\text{Первый мастер}&x&1/x&1\\
\text{Второй мастер}&x+6&1/(x+6)&1\\
\text{Вместе}&4&1/4&1\\
\bottomrule
\end{array}
\]
Складываем \textbf{производительности} мастеров:
\[
\frac1x+\frac1{x+6}=\frac14,\quad x>0.
\]
Умножая на $4x(x+6)$, получаем
\[
4(x+6)+4x=x(x+6),\quad x^2-2x-24=0,
\quad (x-6)(x+4)=0.
\]
Корень $x=-4$ не подходит по смыслу. \textbf{Ответ:} первый мастер --- 6 часов, второй --- 12 часов. Проверка: $1/6+1/12=1/4$.
\secname{5. Смеси: выпаривание и добавление раствора}
Масса растворённого вещества равна произведению массы раствора на массовую долю:
\[
 m_{\text{вещества}}=m_{\text{раствора}}\cdot c,
\qquad c=\frac{p}{100}\quad\text{для концентрации }p\%.
\]
При выпаривании \textbf{только воды} масса соли сохраняется, а масса раствора уменьшается. При добавлении раствора увеличиваются обе массы.

\ex{В 8 кг 30\%-го раствора выпарили 2 кг воды, затем добавили 4 кг 10\%-го раствора. Найдём новую концентрацию.}
\begin{center}
\small
\begin{tabularx}{\linewidth}{>{\raggedright\arraybackslash}Xccc}
\toprule
Этап & Масса раствора, кг & Доля соли & Масса соли, кг \\
\midrule
Исходный раствор & $8$ & $0{,}30$ & $2{,}4$ \\
Испарившаяся вода & $-2$ & $0$ & $0$ \\
Добавленный раствор & $+4$ & $0{,}10$ & $+0{,}4$ \\
\midrule
Итоговая смесь & $10$ & $x$ & $2{,}8$ \\
\bottomrule
\end{tabularx}
\end{center}
Отсюда $10x=8\cdot0{,}30-2\cdot0+4\cdot0{,}10=2{,}8$, значит
\[
\boxed{x=0{,}28=28\%}.
\]
\tip{В условии спрашивают массовую долю соли. Поэтому в конце надо разделить массу соли на массу \textbf{всего} раствора и перевести полученную долю в проценты. Результат «2,8 кг» отвечает на другой вопрос.}

\secname{6. Частые ошибки и самопроверка}
\begin{itemize}
\item $a^ma^n=a^{m+n}$, но $(a^m)^n=a^{mn}$; учитывайте \emph{все} скобки.
\item $(x-3)^2=x^2-6x+9$; не путайте с $x^2-9=(x-3)(x+3)$.
\item Для прямой после нахождения $k,b$ проверьте подстановкой \emph{обе} известные точки.
\item Для вектора $\overrightarrow{AB}$ берите координаты $B-A$, а для скалярного произведения сопоставляйте одноимённые координаты.
\item В задаче о работе сумма производительностей равна производительности совместной работы.
\item В задаче о растворе проверяйте, какая величина требуется в ответе: масса, доля или процент.
\item Не забывайте ограничения: знаменатели не равны нулю, время работы положительно.
\end{itemize}
\clearpage
\thispagestyle{fancy}
\begin{center}{\Large\color{darkaccent} Алгоритм: как составить уравнение по тексту}\end{center}
\vspace{13mm}
\begin{center}
\begin{tikzpicture}[>=Latex,thick, font=\small,
 start/.style={ellipse,draw=accent,line width=.9pt,minimum width=4cm,minimum height=1cm,align=center},
 act/.style={rectangle,rounded corners=2pt,draw=accent,line width=.75pt,text width=5.3cm,minimum height=1.45cm,align=center,inner sep=5pt},
 choice/.style={diamond,draw=accent,aspect=2.9,text width=3.3cm,align=center,inner sep=2pt,minimum height=1.15cm},
 line/.style={->,draw=darkaccent,line width=.8pt}]
\node[start] (s) at (0,0) {Прочитать условие};
\node[choice] (d) at (0,-2.55) {Какой тип задачи?};
\node[act] (w1) at (-4.15,-5.05) {Совместная работа\\[2pt] Принять весь объём за $1$, выбрать времена};
\node[act] (m1) at (4.15,-5.05) {Смеси и растворы\\[2pt] Записать массы растворов и концентрации};
\node[act] (w2) at (-4.15,-8.0) {Вычислить производительности $1/t$ и сложить их};
\node[act] (m2) at (4.15,-8.0) {Найти массу соли $m\cdot c$ и учесть выпаренную воду};
\node[act] (eq) at (0,-11.2) {Составить уравнение баланса и решить его};
\node[act] (check) at (0,-13.65) {Проверить ограничения и смысл полученного результата};
\node[start] (stop) at (0,-15.8) {Записать ответ с единицами};
\draw[line] (s)--(d);
\draw[line] (d.west) -| node[pos=.25,above]{работа} (w1.north);
\draw[line] (d.east) -| node[pos=.25,above]{смеси} (m1.north);
\draw[line] (w1)--(w2);
\draw[line] (m1)--(m2);
\draw[line] (w2.south) |- ([yshift=6mm]eq.north) -- (eq.north);
\draw[line] (m2.south) |- ([yshift=6mm]eq.north);
\draw[line] (eq)--(check);
\draw[line] (check)--(stop);
\end{tikzpicture}
\end{center}
\clearpage
\secname{Тренировочный блок: путь А}
\textbf{10 задач:} решайте по порядку от простых преобразований к текстовым моделям. Сначала выполните задания самостоятельно и только затем проверьте ответы.
\begin{enumerate}[label=\textbf{\arabic*.},itemsep=10pt]
\item Упростите выражение $\displaystyle\frac{(3^2)^3\cdot3^{-4}}{\sqrt9}$.
\item Решите уравнение $\displaystyle 9^{x-1}\sqrt{27}=3^{11/2}$.
\item График $f(x)=kx+b$ проходит через $A(-2;7)$ и $B(2;-1)$. Найдите формулу $f(x)$ и значение $f(-3)$.
\item Прямая $y=kx+b$ проходит через $(-1;1)$ и $(3;9)$. Найдите точку её пересечения с прямой $y=-x+6$.
\item Даны $P(-3;2)$ и $Q(1;5)$. Найдите координаты $\overrightarrow{PQ}$ и его длину.
\item Даны $A(0;1)$, $B(2;2)$, $C(-1;3)$. Найдите $\cos\angle BAC$ и определите величину угла.
\item Первый мастер может выполнить работу на 6 часов быстрее второго. Вместе они выполняют её за 4 часа. Сколько часов потребуется каждому по отдельности?
\item Первый мастер выполняет заказ на 3 часа быстрее второго. Вместе они выполняют его за 2 часа. Определите время каждого мастера.
\item К 8 кг 30\%-го раствора после выпаривания 2 кг воды добавили 4 кг 10\%-го раствора. Какова массовая доля соли в итоговом растворе?
\item В 6 кг 20\%-го раствора выпарили 1 кг воды, затем добавили 5 кг 40\%-го раствора. Найдите концентрацию соли в полученной смеси.
\end{enumerate}
\clearpage
\secname{Ответы к тренировочному блоку}
\begin{center}
\renewcommand{\arraystretch}{1.48}
\begin{tabularx}{\linewidth}{>{\centering\arraybackslash}p{1.4cm}X}
\toprule
\textbf{№} & \textbf{Ответ}\\
\midrule
1 & $3$\\
2 & $x=3$\\
3 & $f(x)=-2x+3$; $f(-3)=9$\\
4 & $(1;5)$\\
5 & $\overrightarrow{PQ}=(4;3)$, $|\overrightarrow{PQ}|=5$\\
6 & $\cos\angle BAC=0$, $\angle BAC=90^\circ$\\
7 & Первый мастер: 6 ч; второй мастер: 12 ч\\
8 & Первый мастер: 3 ч; второй мастер: 6 ч\\
9 & $28\%$\\
10 & $32\%$\\
\bottomrule
\end{tabularx}
\end{center}
\vspace{10mm}
\secname{Итоговая памятка}
\begin{center}
\small
\begin{tabularx}{\linewidth}{>{\raggedright\arraybackslash}p{4.7cm}X}
\toprule
\textbf{Сигнал в задаче} & \textbf{Первое полезное действие}\\
\midrule
Разные степени и корни & Привести к общему основанию.\\
Прямая через две точки & Записать систему на $k$ и $b$.\\
Вектор между двумя точками & Вычесть координаты начала из координат конца.\\
Угол между векторами & Найти скалярное произведение и длины.\\
Совместная работа & Сложить производительности.\\
Выпаривание и добавление & Отдельно учитывать массу раствора и массу соли.\\
\bottomrule
\end{tabularx}
\end{center}
\vfill
\begin{center}\footnotesize\color{muted} Лёвин Артём Александрович \quad\textbullet\quad эксклюзивно для Анастасии Павловой\end{center}
\end{document}
